\section{Introduction (provisional motivation)}
Publicly supported innovation intermediaries can connect firms to complementary expertise, but their provision may also overlap with privately supplied services. This tension was explicit in the United Kingdom in 2026. A House of Commons committee warned that Catapult centres, especially in low-R\&D regions, risk competing with local organisations for funding rather than facilitating knowledge diffusion \citep[paras.~80--84]{UKCommittee2026}. The government subsequently defended Catapults' mixture of core public funding, competitively obtained research funding, and commercial revenue, arguing that this structure responds to demand and limits displacement of private activity \citep{UKGovernment2026}. These are competing policy assessments, not causal estimates of crowding out or crowding in.

A parallel design problem arises in the European Digital Innovation Hubs programme. Hubs are locally accessible while connecting clients to expertise available elsewhere in the network \citep{EuropeanCommissionEDIH}. The network's pricing guidance permits free or reduced-price services but requires a service value or price list to account for the corresponding aid; the published worked examples use market-price benchmarks \citep{EDIHFAQ2026}. The real organisations often combine public and commercial funding and need not be owned by governments. Consequently, these institutions motivate, but do not literally instantiate, our stylised four-entrant game.

The economic issue is not merely the number of supported collaborations. A paid specialist intermediary may select technically capable founders, while a free basic provider broadens participation. Such selection changes the composition of cross-region R\&D matches. Successful new ventures may generate consumer benefits and participant profits while reducing incumbent producer profits, with these effects allocated unevenly across jurisdictions. A local government therefore need not prefer the same provision configuration as a planner evaluating both regions.

We study these mechanisms in a two-region game with two independently deciding local governments and two local private intermediaries. The paper characterizes conditional tariff and selection effects, a regional-versus-national public-entry threshold wedge, and a numerically verified instance of changing provision configurations. It does not claim that the generic mechanisms of selection or local externalities are new, or that government construction of a physical hub is necessarily optimal.

\paragraph{Drafting status.} This is a Stage~9 evidence-linked seed, not the final Stage~10 Introduction. The final literature positioning and exposition order will be determined after writing the model, equilibrium and proof sections.
